% Figure: the three Bernoulli heads (velocity, pressure, elevation) along a
% streamline through a contraction, showing total head constant.
\begin{figure}[htbp]
\centering
\begin{tikzpicture}
\begin{axis}[
  width=12cm, height=7cm,
  xlabel={position along streamline $s$ (arb.\ units)},
  ylabel={head (m of water)},
  xmin=0, xmax=10, ymin=0, ymax=11,
  axis lines=left,
  legend pos=south west,
  legend cell align=left,
  grid=major, grid style={enggray!20},
  tick label style={font=\scriptsize},
  label style={font=\small},
  legend style={font=\scriptsize},
]
% Velocity head u^2/(2g): peaks at the throat near s=5
\addplot[engblue, very thick, samples=200, domain=0:10]
  {1 + 3.2*exp(-((x-5)^2)/2.2)};
\addlegendentry{velocity head $u^2/2g$}
% Pressure head p/(rho g): drops where flow speeds up
\addplot[engred, very thick, samples=200, domain=0:10]
  {7 - 3.2*exp(-((x-5)^2)/2.2)};
\addlegendentry{pressure head $p/\rho g$}
% Total head (constant): sum of the two above (elevation taken as datum)
\addplot[engblack, dashed, very thick, samples=2, domain=0:10] {8};
\addlegendentry{total head $H = u^2/2g + p/\rho g$}
\node[engblack, font=\scriptsize, anchor=south] at (axis cs:5,8.1) {$H$ constant};
\node[engblue, font=\scriptsize, anchor=south] at (axis cs:5,4.3) {throat};
\end{axis}
\end{tikzpicture}
\caption[Bernoulli heads along a streamline]{The velocity and
pressure terms of the head form of Bernoulli's equation along a
streamline that passes through a throat near $s=5$ (elevation taken
as the datum, $z=0$). As the flow accelerates into the throat the
velocity head $u^2/2g$ (blue) rises and the pressure head $p/\rho g$
(red) falls by the same amount, so their sum, the total head $H$
(dashed), stays constant. Energy converts between pressure and
kinetic forms without loss; the inviscid, steady, incompressible
assumptions are what keep the dashed line flat. A real duct would
show $H$ sloping gently downward as friction removes mechanical
energy. Computed from a Gaussian throat model.}
\label{fig:vol03ch11:bernoulli-streamline}
\end{figure}
